\subsection{Fourier Inversion Formula}

For the remainder of this section, we assume that the group G is finite and that K is an algebraically closed field whose characteristic does not divide the order of G, so that the element \(|G|\) of K is not zero.

The algebra K{[}G{]} is semisimple (Maschke's theorem) and finite-dimensional. Denote by \(\widehat{G}\) the set of classes of simple K{[}G{]}-modules. For every \(\lambda\in\widehat{G}\) , choose a linear representation \((\mathrm{V}_{\lambda},\pi_{\lambda})\) of G such that the associated K{[}G{]}-module has class \(\lambda\) . The set \(\widehat{G}\) is finite, and the vector spaces \(V_{\lambda}\) are finite-dimensional (VIII, p.~141, Example). For any \(\lambda\in\widehat{G}\) , denote by \(d_{\lambda}\) the degree of the representation \(\pi_{\lambda}\) , that is, the dimension of the K-vector space \(V_{\lambda}\) , and denote its character by \(\chi_{\lambda}\) .

Denote by \(\mathrm{F}(\widehat{\mathrm{G}})\) the product algebra \(\prod_{\lambda\in\widehat{G}}\operatorname{End}_{\mathrm{K}}(\mathrm{V}_{\lambda})\) and by \(\overline{F}\) the mapping from K{[}G{]} to \(\mathrm{F}(\widehat{\mathrm{G}})\) defined by \(\overline{\mathcal{F}}(a)=(\pi_{\lambda}(a))_{\lambda\in\widehat{G}}\) . Since the field K is algebraically closed, the mapping \(\overline{F}\) is an algebra isomorphism (loc. cit.).

For every \(\lambda \in \widehat{\mathbf{G}}\) , the dimension of the algebra \(\mathrm{End}_{\mathrm{K}}(\mathrm{V}_{\lambda})\) is \(d_{\lambda}^{2}\) ; that of the algebra \(\mathrm{K}[\mathrm{G}]\) is \(\mathrm{Card}(\mathrm{G})\) . We therefore have the relation

\[
\mathrm{Card(G)} = \sum_ {\lambda \in \widehat {\mathrm{G}}} d _ {\lambda} ^ {2}.\tag{8}
\]

Denote by \(\tau\) the trace in the algebra K{[}G{]}; by definition, the trace \(\tau(a)\) of an element \(a\) of K{[}G{]} is the trace of the endomorphism \(x \mapsto ax\) of K{[}G{]} (III, §7, No.~7, p.~515). Let \(a = \sum_{g \in G} a_g g\) be an element of K{[}G{]}; by formula (2), we have \(\tau(ag^{-1}) = |G| a_g\) for every \(g \in G\) , and therefore the relation

\[
a = | \mathrm{G} | ^ {- 1} \sum_ {g \in \mathrm{G}} \tau (a g ^ {- 1}) g.\tag{9}
\]

Denote by \(\widehat{\tau}\) the trace in the algebra \(F(\widehat{G})\) . Let \(A = (A_{\lambda})_{\lambda \in \widehat{G}}\) be an element of \(F(\widehat{G})\) ; we have (cf.~III, §9, No.~3, p.~545, Example 3)

\[
\widehat {\tau} (\mathrm{A}) = \sum_ {\lambda \in \widehat {\mathrm{G}}} d _ {\lambda} \operatorname{Tr} (\mathrm{A} _ {\lambda}).\tag{10}
\]

Since the mapping \(\overline{F}\) is a K-algebra isomorphism, we have \(\widehat{\tau}\circ\overline{F}=\tau\) , and therefore

\[
\tau (a) = \widehat {\tau} (\overline {{\mathcal {F}}} (a)) = \widehat {\tau} \big ((\pi_ {\lambda} (a)) _ {\lambda \in \widehat {\mathrm{G}}} \big) = \sum_ {\lambda \in \widehat {\mathrm{G}}} d _ {\lambda} \operatorname{Tr} (\pi_ {\lambda} (a)) = \sum_ {\lambda \in \widehat {\mathrm{G}}} d _ {\lambda} \operatorname{Tr} _ {\lambda} (a)
\]

for every \(a\in \mathrm{K}[\mathrm{G}]\) , that is

\[
\tau = \sum_ {\lambda \in \widehat {\mathrm{G}}} d _ {\lambda} \operatorname{Tr} _ {\lambda}.\tag{11}
\]

Therefore, by (2) of VIII, p.~399, for \(g \in \mathbf{G}\) , we have

\[
\sum_ {\lambda \in \widehat {G}} d _ {\lambda}   \chi_ {\lambda} (g) = \left\{ \begin{array}{l l} | G | & \text { if   } g \text {   is   the   identity   element }, \\ 0 & \text { otherwise }. \end{array} \right.\tag{12}
\]

For \(a \in \mathrm{K}[\mathrm{G}]\) , relation (9) takes on the following form:

\[
a = | \mathrm{G} | ^ {- 1} \sum_ {g \in \mathrm{G}} \sum_ {\lambda \in \widehat {\mathrm{G}}} d _ {\lambda} \left. \operatorname{Tr} (\pi_ {\lambda} (a) \pi_ {\lambda} (g ^ {- 1})) g; \right.\tag{13}
\]

it follows that for every element \(A = (A_{\lambda})_{\lambda \in \widehat{G}}\) of \(F(\widehat{G})\) , we have

\[
\overline {{\mathcal {F}}} ^ {- 1} (\mathrm{A}) = | \mathrm{G} | ^ {- 1} \sum_ {g \in \mathrm{G}} \sum_ {\lambda \in \widehat {\mathrm{G}}} d _ {\lambda} \operatorname{Tr} (\mathrm{A} _ {\lambda} \pi_ {\lambda} (g ^ {- 1})) g\tag{14}
\]

(``Fourier inversion formula'').

For \(\mu \in \widehat{\mathrm{G}}\) , denote by \(j_{\mu}:\mathrm{End}_{\mathrm{K}}(\mathrm{V}_{\mu})\longrightarrow \prod_{\lambda \in \widehat{\mathrm{G}}} \mathrm{End}_{\mathrm{K}}(\mathrm{V}_{\lambda})\) the mapping such that \(j_{\mu}(u) = (v_{\lambda})\) , where \(v_{\lambda} = 0\) if \(\lambda \neq \mu\) and \(v_{\mu} = u\) . By formula (14), we have

\[
\overline {{\mathcal {F}}} ^ {- 1} (j _ {\mu} (u)) = | \mathrm{G} | ^ {- 1} d _ {\mu} \sum_ {g \in \mathrm{G}} \mathrm{Tr} (u \pi_ {\mu} (g ^ {- 1})) g.\tag{15}
\]

The center of the algebra \(\mathrm{F}(\widehat{\mathrm{G}})=\prod_{\lambda\in\widehat{\mathrm{G}}} \mathrm{End}_{\mathrm{K}}(\mathrm{V}_{\lambda})\) consists of the families \((a_{\lambda}1_{\mathrm{V}_{\lambda}})_{\lambda\in\widehat{\mathrm{G}}}\) , where \((a_{\lambda})\) is a family of elements of K. It is the image by \(\overline{{\mathcal{F}}}\) of the center of the algebra K{[}G{]}. That center therefore has a basis \((e_{\lambda})_{\lambda\in\widehat{\mathrm{G}}}\) characterized by the relation

\[
\pi_ {\lambda} (e _ {\mu}) = \delta_ {\lambda \mu} 1 _ {V _ {\lambda}}\tag{16}
\]

for \(\lambda, \mu \in \widehat{\mathbf{G}}\) , where \(\delta_{\lambda \mu}\) is the Kronecker delta function. By formula (15), we have, for every \(\mu \in \widehat{\mathbf{G}}\) ,

\[
e _ {\mu} = | \mathrm{G} | ^ {- 1} d _ {\mu} \sum_ {g \in \mathrm{G}} \chi_ {\mu} (g ^ {- 1}) g.\tag{17}
\]

These elements satisfy the relations

\[
\sum_ {\lambda \in \widehat {\mathsf {G}}} e _ {\lambda} = 1, \quad e _ {\mu} ^ {2} = e _ {\mu}, \quad \mathrm{and} \quad e _ {\mu} e _ {\nu} = 0\tag{18}
\]

for all \(\mu, \nu \in \widehat{\mathrm{G}}\) such that \(\mu \neq \nu\) ; they are the indecomposable idempotents of \(\mathrm{Z}(\mathrm{K}[\mathrm{G}])\) (VIII, p.~147, Proposition 15)

Remark. --- Let \((V, \pi)\) be a linear representation of G. By Maschke's theorem, the K{[}G{]}-module V is semisimple. For any \(\lambda \in \widehat{G}\) , denote by \(V^{\lambda}\) the isotypical component of type \(\lambda\) of the K{[}G{]}-module V; we have \(V = \bigoplus_{\lambda \in \widehat{G}} V^{\lambda}\) . By Proposition 15 of VIII, p.~147 and formula (17), the projector of V with image \(V^{\lambda}\) associated with this decomposition of V is equal to

\[
\pi (e _ {\lambda}) = | \mathrm{G} | ^ {- 1} d _ {\lambda} \sum_ {g \in \mathrm{G}} \chi_ {\lambda} (g ^ {- 1}) \pi (g).\tag{19}
\]

Applying this formula to \(\pi = \pi_{\lambda}\) , we obtain that the element \(d_{\lambda} \cdot 1\) of K is not zero. We will see further on (VIII, p.~420, Corollary 2) that \(d_{\mu}\) divides the cardinal of G.

Let \(\lambda\) be an element of \(\widehat{\mathbf{G}}\) . The mapping \(\pi_{\lambda}\) from K{[}G{]} to \(\mathrm{End}_{\mathrm{K}}(\mathrm{V}_{\lambda})\) induces an isomorphism of (K{[}G{]}, K{[}G{]})-bimodules from \(e_{\lambda}\mathrm{K}[\mathrm{G}]\) to \(\mathrm{End}_{\mathrm{K}}(\mathrm{V}_{\lambda})\) . By Proposition 15 of VIII, p.~147, the K{[}G{]}-submodule \(e_{\lambda}\mathrm{K}[\mathrm{G}]\) is the isotypical component of K{[}G{]} of type \(\lambda\) . It is also the isotypical component of type \(\lambda^{\vee}\) for the right regular representation of G (VIII, p.~143, Proposition 11) as well as the isotypical component of type \(\lambda \times \lambda^{\vee}\) for the biregular representation (VIII, p.~381, Proposition 4).

